\chapter{Dérivation de la matrice $T$ à partir de la théorie de la perturbation à plusieurs corps}\label{annexe:défauts}

Cette annexe présente la dérivation complète de la matrice $T$ introduite à la section~\ref{sec:mbt_défaut}. Nous partons de l'Hamiltonien de couplage entre les électrons et un seul défaut, construisons la fonction de Green du système perturbée par resommation de la série de Dyson, puis généralisons le problème à un ensemble dilué de défauts afin de simuler un matériau désordonné. 

\section{Hamiltonien et fonction de Green pour un seul défaut}
Conformément à la section~\ref{sec:hamiltonien}, l'Hamiltonien d'électrons en présence d'un seul défaut se décompose en une partie libre et une partie de couplage
\begin{equation}
    \hat{H} = \hat{H}_0+\hat{H}_\text{ed} = \sum_{n}\varepsilon_{n}\hat{c}^\dag_{n}\hat{c}_{n}+\sum_{mn}M_{mn}\hat{c}^\dag_{m}\hat{c}_{n},
\end{equation}
où $\varepsilon_{n}$ sont les énergies propres de Kohn-Sham du cristal parfait et où les éléments de matrice électron-défaut $M_{mn}$ contiennent les contributions locale et non-locale dérivées au chapitre~\ref{chap:défauts}. Nous avons également supprimé les indices de vecteur d'onde pour alléger la notation dans cette annexe. Nous allons utiliser la théorie de la perturbation en mécanique quantique sur l'Hamiltonien de couplage pour résoudre ce problème.

Le lecteur peut se questionner sur l'utilité d'employer la théorie de perturbation pour ce problème, qui peut être résolu exactement. En effet, l'Hamiltonien de couplage est quadratique en les opérateurs de création et d'annihilation électroniques. Cela est une conséquence du fait qu'en présence d'un seul défaut, le problème est essentiellement un problème à un corps. Cependant, lorsque nous ajouterons plusieurs défauts, le problème ne peut plus être résolu de manière exacte et la théorie de la perturbation nous permet de résoudre le problème approximativement. La quantité centrale dans la théorie de la perturbation à plusieurs corps est la fonction de Green de Matsubara du système perturbé
\begin{equation}
    \mathcal{G}_{mn}(\tau'-\tau) = -\frac{{}_{0}\langle\mathcal{T}\hat{c}_m(\tau')S
    (\beta, 0)\hat{c}^\dag_n(\tau)\rangle_0}{{}_{0}\langle S(\beta, 0) \rangle_0},\label{eqn:G}
\end{equation}
où ${}_{0}\langle\dots\rangle_0$ est une moyenne thermique par rapport à l'Hamiltonien non-perturbé $\hat{H}_0$, $\tau$ est un temps imaginaire, $\mathcal{T}$ est l'opérateur d'ordonnancement temporel (il ordonne les opérateurs dans le temps, avec le plus grand temps à gauche), $\beta=1/k_BT$ est l'inverse de la température, et la matrice $S$ est définie formellement comme étant
\begin{equation}
    S(\beta, 0) = \mathcal{T}\exp{-\int_0^\beta d\tau\, \hat{H}_\text{ed}(\tau)} = \sum_{p=0}^\infty \frac{(-1)^p}{p!}\int_0^\beta d\tau_1\dots\int_0^\beta d\tau_p\,\hat{H}_\text{ed}(\tau_1)\dots\hat{H}_\text{ed}(\tau_p),
\end{equation}
où les opérateurs $\hat{O}(\tau)$ sont écrits dans la représentation d'interaction. Pour évaluer la fonction de Green, nous développons la matrice $S$ en puissances de l'interaction $\hat{H}_\text{ed}$, utilisons le théorème de Wick pour apparier les opérateurs électroniques, puis resommons la série ainsi engendrée à l'aide de diagrammes de Feynman.

\section{Équation de Dyson}
Avant d'écrire la série perturbative pour la fonction de Green de Matsubara, il est utile de formuler quelques règles pour évaluer correctement les termes de cette série.
\begin{itemize}
    \item Premièrement, nous allons apparier les opérateurs électroniques en paires. Comme ces paires vont être moyennées par rapport à l'état fondamental, les seules paires non-nulles sont celles qui contiennent exactement un opérateur de création et un opérateur d'annihilation. 
    \item De plus, pour une paire donnée ${}_{0}\langle \mathcal{T}\hat{c}_m(\tau')\hat{c}_n^\dag(\tau)\rangle_0$, l'effet de l'opérateur $\hat{c}^\dag_n(\tau)$ est de placer un électron dans la bande $n$. Pour que cette moyenne soit non-nulle, il est impératif que l'opérateur d'annihilation enlève cet électron, ce qui impose $m=n$. Ainsi,
    \[{}_{0}\langle \mathcal{T}\hat{c}_m(\tau')\hat{c}_n^\dag(\tau)\rangle_0=\delta_{mn}\, {}_{0}\langle \mathcal{T}\hat{c}_n(\tau')\hat{c}_n^\dag(\tau)\rangle_0.\]
    \item Les deux prochaines règles sont à propos de l'ordonnancement temporel des opérateurs. Premièrement, pour la moyenne d'une paire évaluée à temps égal, il est de coutume de placer l'opérateur de création à gauche, donnant la distribution de Fermi-Dirac pour des opérateurs fermioniques comme les électrons
    \[{}_0\langle\mathcal{T}\hat{c}_m(\tau)\hat{c}^\dag_n(\tau)\rangle_0=-\delta_{mn}f_n.\]
    \item Deuxièmement, pour la moyenne d'une paire évaluée à des temps différents, il est de coutume de placer l'opérateur de création à droite, donnant la fonction de Green de Matsubara pour le système non-perturbé
    \[{}_0\langle \mathcal{T}\hat{c}_m(\tau')\hat{c}^\dag_n(\tau)\rangle_0=-\delta_{mn} \mathcal{G}_n^{(0)}(\tau'-\tau).\]
    Nous notons qu'il faut ajouter un signe moins pour chaque permutation d'opérateurs fermioniques dans les deux cas, et qu'il y a un autre signe moins dans la définition de la fonction de Green de Matsubara.
    \item Finalement, lorsque nous apparions des opérateurs fermioniques avec le théorème de Wick, il faut ajouter un signe moins pour chaque permutation nécessaire pour mettre les opérateurs les uns à côté des autres. Cela devient rapidement lourd et il est très facile de faire des erreurs. Une façon plus visuelle est de compter le nombre de croisements en reliant deux opérateurs par une ligne.
\end{itemize} 

Voyons comment écrire la série pour la fonction de Green de Matsubara pour le système perturbé à l'aide des règles décrites ci-haut. En principe, nous devons faire l'expansion de la matrice $S$ au numérateur et au dénominateur simultanément, ce qui devient intraitable très rapidement. Heureusement, en vertu du théorème des amas connexes (traduit de l'anglais \textit{linked-cluster theorem}), les diagrammes engendrés par l'expansion au dénominateur annulent exactement les diagrammes déconnectés engendrés par l'expansion du numérateur. Ainsi, il suffit de ne considérer que l'expansion du numérateur et de ne garder que les diagrammes connectés. Commençons par le numérateur de l'équation~\eqref{eqn:G} et procédons ordre par ordre. 

À l'ordre zéro, nous n'avons pas de puissance de l'interaction et ce terme est simplement
\begin{align}
    \mathcal{G}_{mn}^{(0)} &= -{}_0\langle\mathcal{T}\hat{c}_m(\tau')\hat{c}^\dag_n(\tau)\rangle_0,\\
    &= -\delta_{mn}\,{}_0\langle\mathcal{T}\hat{c}_n(\tau')\hat{c}^\dag_n(\tau)\rangle_0,\\
    &= \delta_{mn}\mathcal{G}^{(0)}_n(\tau'-\tau),
\end{align}
soit la fonction de Green de Matsubara du système non-perturbé. Pour tracer le diagramme de Feynman qui représente ce terme, nous utilisons les règles de Feynman développées pour la spectroscopie Raman au chapitre~\ref{chap:Raman}. Le diagramme correspondant est illustré à la figure~\ref{fig:dyson_ordre0}.
\begin{figure}[h]
    \centering
    % Diagramme de Feynman : terme d'ordre zéro de la série de Dyson
% pour la fonction de Green électronique en présence d'un défaut.
% Style identique aux diagrammes du chapitre Raman (tikz/e1.tex, tikz/pd.tex).
% Compatible pdflatex. Requiert : tikz, decorations.markings, arrows.meta, calc.
\begin{tikzpicture}[
    line/.style={thick},
    >=Stealth,
    fermion/.style={
      postaction={decorate},
      decoration={markings, mark=at position 0.55 with {\arrow{Stealth}}}
    },
]
  % --- propagateur libre ---
  \draw[line, fermion] (0,0) -- (4,0);
  \node[above=3pt] at (2,0) {$n$};

  % --- extrémités ---
  \node[below=4pt] at (0,0) {$\tau$};
  \node[below=4pt] at (4,0) {$\tau'$};
  \fill (0,0) circle (0.05);
  \fill (4,0) circle (0.05);
\end{tikzpicture}

    \caption[Diagramme de Feynman du terme d'ordre zéro de la série de Dyson.]{Diagramme de Feynman du terme d'ordre zéro de la série de Dyson. L'électron se propage librement dans la bande $n$, du temps $\tau$ au temps $\tau'$, sans interagir avec le défaut.}
    \label{fig:dyson_ordre0}
\end{figure}
\noindent
Physiquement, ce diagramme représente un électron dans la bande $n$ au temps $\tau$ et qui évolue dans le temps vers $\tau'$ tout en restant dans la bande $n$. C'est donc un électron qui se propage librement sans interagir avec le défaut. Ce diagramme est connecté.

Le terme d'ordre un s'écrit
\begin{align}
    \mathcal{G}_{mn}^{(1)}(\tau'-\tau) &= (-1)^2\int_0^\beta d\tau_1 \,{}_0\langle\mathcal{T}\hat{c}_m(\tau')\hat{H}_\text{ed}(\tau_1)\hat{c}^\dag_n(\tau)\rangle_0,\\
    &= (-1)^2\sum_{n_1m_1}\int_0^\beta d\tau_1\,{}_0\langle\mathcal{T}\hat{c}_m(\tau')\hat{c}^\dag_{m_1}(\tau_1)\hat{c}_{n_1}(\tau_1)\hat{c}^\dag_n(\tau)\rangle_0.
\end{align}
En utilisant le théorème de Wick, il n'y a que deux appariements qui sont non-nuls. Ils sont représentés graphiquement par
\begin{equation}
\begin{aligned}
\mathcal{G}_{mn}^{(1)}(\tau'-\tau)
= (-1)^2 \sum_{m_1n_1} \int_0^\beta &d\tau_1\, 
\Big[
    {}_0\langle\mathcal{T}
      \wick{
        \c1 c_{m}(\tau')
        \c1 c_{m_1}^\dag(\tau_1)
        \c2 c_{n_1}(\tau_1)
        \c2 c_{n}^\dag(\tau)
      }\rangle_0
\\
&
    +
    {}_0\langle\mathcal{T}
    \wick{
        \c1 c_{m}(\tau')
        \c2 c_{m_1}^\dag(\tau_1)
        \c2 c_{n_1}(\tau_1)
        \c1 c_{n}^\dag(\tau)
    }\rangle_0
\Big]M_{m_1n_1}.
\end{aligned}
\end{equation}
Ici, les lignes de contraction de Wick ne se croisent pas pour le premier terme et se croisent deux fois pour le deuxième terme, il faut donc ajouter un nombre pair de signes moins au deuxième terme, ce qui ne change rien. Les appariements écrits explicitement sont
\begin{equation}
\begin{aligned}
\mathcal{G}_{mn}^{(1)}(\tau'-\tau)
= (-1)^2 \sum_{m_1n_1} \int_0^\beta &d\tau_1\, 
\Big[
    {}_0\langle\mathcal{T}\hat{c}_{m}(\tau')\hat{c}_{m_1}^\dag(\tau_1)
      \rangle_0    
      {}_0\langle\mathcal{T}\hat{c}_{n_1}(\tau_1)\hat{c}_{n}^\dag(\tau)
      \rangle_0
\\
&
    +
    {}_0\langle\mathcal{T}\hat{c}_{m}(\tau')\hat{c}_{n}^\dag(\tau)
      \rangle_0    
      {}_0\langle\mathcal{T}\hat{c}_{m_1}(\tau_1)\hat{c}_{n_1}^\dag(\tau_1)
      \rangle_0
\Big]M_{m_1n_1}.
\end{aligned}
\end{equation}
Les moyennes thermiques des paires d'opérateurs sont ensuite évaluées selon les règles explicitées plus haut. Cela donne
\begin{equation}
\begin{aligned}
\mathcal{G}_{mn}^{(1)}(\tau'-\tau)
= \int_0^\beta &d\tau_1\, 
\Big[
    \mathcal{G}^{(0)}_m(\tau'-\tau_1)\mathcal{G}^{(0)}_{n}(\tau_1-\tau)M_{mn}
\\
&
    -\delta_{mn}\sum_{n_1}\mathcal{G}_{n}^{(0)}(\tau'-\tau)f_{n_1} M_{n_1n_1}
\Big].
\end{aligned}
\end{equation}
Il y a donc deux termes qui contribuent au premier ordre et ils peuvent être représentés par la somme des deux diagrammes de Feynman de la figure~\ref{fig:dyson_ordre1}.
\begin{figure}[h]
    \centering
    % Diagrammes de Feynman : termes du premier ordre de la série de Dyson
% pour la fonction de Green électronique en présence d'un défaut.
% (a) diagramme connecté : une diffusion sur le défaut au temps tau_1.
% (b) diagramme déconnecté (tadpole) : boucle fermée au temps tau_1.
% Style identique aux diagrammes du chapitre Raman (tikz/e1.tex, tikz/pd.tex).
% Compatible pdflatex. Requiert : tikz, decorations.markings, arrows.meta, calc.

% Vertex électron-défaut : cercle blanc + croix (cf. tikz/règles_Raman.tex)
\providecommand{\vertexdefaut}[1]{%
  \draw[fill=white] (#1) circle (0.2);
  \draw[line width=1pt] ($(#1)+(-0.14,-0.14)$) -- ($(#1)+(0.14,0.14)$);
  \draw[line width=1pt] ($(#1)+(-0.14,0.14)$) -- ($(#1)+(0.14,-0.14)$);
}

\begin{tikzpicture}[
    line/.style={thick},
    >=Stealth,
    fermion/.style={
      postaction={decorate},
      decoration={markings, mark=at position 0.55 with {\arrow{Stealth}}}
    },
]

% ===================== (a) diagramme connecté =====================
\begin{scope}[shift={(0,0)}]
  \draw[line, fermion] (0,0)   -- (1.82,0);
  \draw[line, fermion] (2.18,0) -- (4,0);
  \node[above=3pt] at (0.9,0) {$n$};
  \node[above=3pt] at (3.1,0) {$m$};

  \vertexdefaut{2,0}
  \node[above=6pt] at (2,0) {$\tau_1$};
  \node[below=6pt] at (2,0) {$M_{mn}$};

  \node[below=4pt] at (0,0) {$\tau$};
  \node[below=4pt] at (4,0) {$\tau'$};
  \fill (0,0) circle (0.05);
  \fill (4,0) circle (0.05);

  \node[font=\small\bfseries] at (-0.6,1.6) {(a)};
\end{scope}

% ===================== (b) diagramme déconnecté =====================
\begin{scope}[shift={(7.2,0)}]
  \draw[line, fermion] (0,0) -- (4,0);
  \node[below=3pt] at (1.2,0) {$n$};

  % boucle fermée, centre (2,1.15), rayon 0.55
  \draw[line, fermion] (2,0.60) arc (-90:90:0.55);
  \draw[line, fermion] (2,1.70) arc (90:270:0.55);
  \node[left=3pt] at (1.45,1.15) {$n_1$};

  \vertexdefaut{2,1.70}
  \node[above=6pt] at (2,1.70) {$\tau_1$, $M_{n_1n_1}$};

  \node[below=4pt] at (0,0) {$\tau$};
  \node[below=4pt] at (4,0) {$\tau'$};
  \fill (0,0) circle (0.05);
  \fill (4,0) circle (0.05);

  \node[font=\small\bfseries] at (-0.6,1.6) {(b)};
\end{scope}

\end{tikzpicture}

    \caption[Diagrammes de Feynman des termes du premier ordre de la série de Dyson.]{Diagrammes de Feynman des termes du premier ordre de la série de Dyson. (a) Diagramme connecté : l'électron diffuse une fois sur le défaut au temps $\tau_1$, avec une amplitude $M_{mn}$. (b) Diagramme déconnecté : la boucle fermée au temps $\tau_1$ n'est reliée par aucun propagateur au reste du diagramme et est annulée par le dénominateur de l'équation~\eqref{eqn:G}.}
    \label{fig:dyson_ordre1}
\end{figure}
\noindent
Physiquement, le premier diagramme représente un électron dans la bande $n$ au temps $\tau$ qui se propage vers la bande $m$ au temps $\tau'$ en diffusant sur un défaut au temps $\tau_1$ avec une amplitude de probabilité $M_{mn}$. Le deuxième diagramme est déconnecté et nous pouvons l'ignorer.

Au deuxième ordre, l'expansion du numérateur s'écrit
\begin{align}
    \mathcal{G}^{(2)}_{mn}(\tau'-\tau) &= \frac{(-1)^3}{2!}\int_0^\beta d\tau_1\,\int_0^\beta d\tau_2\, {}_0\langle\mathcal{T}\hat{c}_m(\tau')\hat{H}_\text{ed}(\tau_1)\hat{H}_\text{ed}(\tau_2)\hat{c}^\dag_{n}(\tau)\rangle_0\\
    & = \frac{(-1)^3}{2!}\sum_{n_1m_1}\sum_{n_2m_2}\int_0^\beta d\tau_1\,\int_0^\beta d\tau_2\, {}_0\langle\mathcal{T}\hat{c}_m(\tau')\hat{c}_{m_1}^\dag(\tau_1)\hat{c}_{n_1}(\tau_1)\hat{c}_{m_2}^\dag(\tau_2)\hat{c}_{n_2}(\tau_2)\hat{c}_n^\dag(\tau)\rangle_0M_{m_1n_1}M_{m_2n_2}.
\end{align}
Il y a maintenant six appariements non-nuls. En général, pour $N$ paires d'opérateurs de création et d'annihilation fermioniques, il y a $N!$ appariements non-nuls. Ces appariements sont
\begin{equation}
\begin{aligned}
    \mathcal{G}_{mn}^{(2)}(\tau'-\tau) &= \frac{(-1)^3}{2!}\sum_{n_1m_1}\sum_{n_2m_2}\int_0^\beta d\tau_1\,\int_0^\beta d\tau_2\,\Big[\\
    &\wick{{}_0\langle\mathcal{T} \c1 c_{m}(\tau') \c1 c_{m_1}^\dag(\tau_1) \c1 c_{n_1}(\tau_1) \c1 c_{m_2}^\dag(\tau_2) \c1 c_{n_2}(\tau_2) \c1 c_{n}^\dag(\tau)}\rangle_0\\
    &+\wick{{}_0\langle\mathcal{T} \c1 c_{m}(\tau') \c1 c_{m_1}^\dag(\tau_1) \c1 c_{n_1}(\tau_1) \c2 c_{m_2}^\dag(\tau_2) \c2 c_{n_2}(\tau_2)\c1 c_{n}^\dag(\tau)}\rangle_0\\
    &+\wick{{}_0\langle\mathcal{T} \c1 c_{m}(\tau') \c2 c_{m_1}^\dag(\tau_1) \c2 c_{n_1}(\tau_1) \c1 c_{m_2}^\dag(\tau_2) \c1 c_{n_2}(\tau_2) \c1 c_{n}^\dag(\tau)}\rangle_0\\
    &+\wick{{}_0\langle\mathcal{T} \c1 c_{m}(\tau') \c2 c_{m_1}^\dag(\tau_1) \c3 c_{n_1}(\tau_1) \c1 c_{m_2}^\dag(\tau_2) \c2 c_{n_2}(\tau_2) \c3 c_{n}^\dag(\tau)}\rangle_0\\
    &+\wick{{}_0\langle\mathcal{T} \c1 c_{m}(\tau')\c2 c_{m_1}^\dag(\tau_1)\c2 c_{n_1}(\tau_1)\c2 c_{m_2}^\dag(\tau_2)\c2 c_{n_2}(\tau_2) \c1 c_{n}^\dag(\tau)}\rangle_0\\
    &+\wick{{}_0\langle\mathcal{T} \c1 c_{m}(\tau') \c2 c_{m_1}^\dag(\tau_1)\c3 c_{n_1}(\tau_1)\c3 c_{m_2}^\dag(\tau_2)\c2 c_{n_2}(\tau_2)\c1 c_{n}^\dag(\tau)}\rangle_0\Big]M_{m_1n_1}M_{m_2n_2}.
\end{aligned}
\end{equation}
Ici, tous les termes sauf le quatrième ont un nombre pair de croisements ; nous devons donc ajouter un signe moins au quatrième terme uniquement. En écrivant les paires d'opérateurs fermioniques explicitement, nous avons
\begin{equation}
\begin{aligned}
    \mathcal{G}^{(2)}(\tau'-\tau) &= \frac{(-1)^3}{2!}\sum_{n_1m_1}\sum_{n_2m_2}\int_0^\beta d\tau_1\,\int_0^\beta d\tau_2\Big[\\
    &{}_0\langle\mathcal{T}\hat{c}_{m}(\tau')\hat{c}_{m_1}^\dag(\tau_1)\rangle_0{}_0\langle\mathcal{T}\hat{c}_{n_1}(\tau_1)\hat{c}_{m_2}^\dag(\tau_2)\rangle_0{}_0\langle\mathcal{T}\hat{c}_{n_2}(\tau_2)\hat{c}_{n}^\dag(\tau)\rangle_0\\
    &+{}_0\langle\mathcal{T}\hat{c}_{m}(\tau')\hat{c}_{m_1}^\dag(\tau_1)\rangle_0{}_0\langle\mathcal{T}\hat{c}_{n_1}(\tau_1)\hat{c}_{n}^\dag(\tau)\rangle_0{}_0\langle\mathcal{T}\hat{c}_{m_2}^\dag(\tau_2)\hat{c}_{n_2}(\tau_2)\rangle_0\\
    &+{}_0\langle\mathcal{T}\hat{c}_{m}(\tau')\hat{c}_{m_2}^\dag(\tau_2)\rangle_0{}_0\langle\mathcal{T}\hat{c}_{m_1}^\dag(\tau_1)\hat{c}_{n_1}(\tau_1)\rangle_0{}_0\langle\mathcal{T}\hat{c}_{n_2}(\tau_2)\hat{c}_{n}^\dag(\tau)\rangle_0\\
    &+{}_0\langle\mathcal{T}\hat{c}_{m}(\tau')\hat{c}_{m_2}^\dag(\tau_2)\rangle_0{}_0\langle\mathcal{T}\hat{c}_{m_1}^\dag(\tau_1)\hat{c}_{n_2}(\tau_2)\rangle_0{}_0\langle\mathcal{T}\hat{c}_{n_1}(\tau_1)\hat{c}_{n}^\dag(\tau)\rangle_0\\
    &+{}_0\langle\mathcal{T}\hat{c}_{m}(\tau')\hat{c}_{n}^\dag(\tau)\rangle_0{}_0\langle\mathcal{T}\hat{c}_{m_1}^\dag(\tau_1)\hat{c}_{n_1}(\tau_1)\rangle_0{}_0\langle\mathcal{T}\hat{c}_{m_2}^\dag(\tau_2)\hat{c}_{n_2}(\tau_2)\rangle_0\\
    &+{}_0\langle\mathcal{T}\hat{c}_{m}(\tau')\hat{c}_{n}^\dag(\tau)\rangle_0{}_0\langle\mathcal{T}\hat{c}_{m_1}^\dag(\tau_1)\hat{c}_{n_2}(\tau_2)\rangle_0{}_0\langle\mathcal{T}\hat{c}_{n_1}(\tau_1)\hat{c}_{m_2}^\dag(\tau_2)\rangle_0\Big]\\
    &\times M_{m_1n_1}M_{m_2n_2}.
\end{aligned}
\end{equation}
Finalement, en évaluant les moyennes thermiques nous obtenons
\begin{equation}
\begin{aligned}
    \mathcal{G}^{(2)}_{mn}(\tau'-\tau) &= \frac{1}{2}\int_0^\beta d\tau_1\,\int_0^\beta d\tau_2\Big[\\
    &\mathcal{G}^{(0)}_m(\tau'-\tau_1)M_{mn_1}\mathcal{G}^{(0)}_{n_1}(\tau_1-\tau_2)M_{n_1n}\mathcal{G}_n^{(0)}(\tau_2-\tau)\\
    & -\sum_{n_2}\mathcal{G}^{(0)}_m(\tau'-\tau_1)M_{mn}\mathcal{G}^{(0)}_n(\tau_1-\tau)f_{n_2}M_{n_2n_2}\\
    &-\sum_{n_1}\mathcal{G}^{(0)}_m(\tau'-\tau_2)M_{mn}\mathcal{G}^{(0)}_n(\tau_2-\tau)f_{n_1}M_{n_1n_1}\\
    &+\mathcal{G}^{(0)}_m(\tau'-\tau_2)M_{mn_2}\mathcal{G}_{n_2}^{(0)}(\tau_2-\tau_1)M_{n_2n}\mathcal{G}_n^{(0)}(\tau_1-\tau)\\
    &+\delta_{mn}\sum_{n_1n_2}\mathcal{G}_n^{(0)}(\tau'-\tau)f_{n_1}M_{n_1n_1}f_{n_2}M_{n_2n_2}\\
    &-\delta_{mn}\sum_{n_1n_2}\mathcal{G}^{(0)}_n(\tau'-\tau)M_{n_1n_2}\mathcal{G}_{n_1}^{(0)}(\tau_1-\tau_2)M_{n_2n_1}\mathcal{G}^{(0)}_{n_2}(\tau_2-\tau_1)
    \Big].
\end{aligned}
\end{equation}
Ces six appariements se traduisent par les six diagrammes de Feynman de la figure~\ref{fig:dyson_ordre2}.
\begin{figure}[h]
    \centering
    % Diagrammes de Feynman : les six termes du deuxième ordre de la série de Dyson
% pour la fonction de Green électronique en présence d'un défaut.
% L'ordre des panneaux suit celui des six appariements de Wick de l'annexe.
%   (a) double diffusion, tau < tau_2 < tau_1 < tau'
%   (b) une diffusion au temps tau_1 + boucle fermée au temps tau_2
%   (c) une diffusion au temps tau_2 + boucle fermée au temps tau_1
%   (d) double diffusion, tau < tau_1 < tau_2 < tau'
%   (e) propagation libre + deux boucles fermées
%   (f) propagation libre + boucle à deux vertex
% Seuls (a) et (d) sont connectés ; les autres sont annulés par le dénominateur.
% Style identique aux diagrammes du chapitre Raman (tikz/e1.tex, tikz/pd.tex).
% Compatible pdflatex. Requiert : tikz, decorations.markings, arrows.meta, calc.

% Vertex électron-défaut : cercle blanc + croix (cf. tikz/règles_Raman.tex)
\providecommand{\vertexdefaut}[1]{%
  \draw[fill=white] (#1) circle (0.2);
  \draw[line width=1pt] ($(#1)+(-0.14,-0.14)$) -- ($(#1)+(0.14,0.14)$);
  \draw[line width=1pt] ($(#1)+(-0.14,0.14)$) -- ($(#1)+(0.14,-0.14)$);
}

% Boucle fermée (tadpole) : #1 centre x, #2 centre y, #3 rayon, #4 indice de bande
\providecommand{\boucledefaut}[4]{%
  \draw[line, fermion] (#1,{#2-#3}) arc (-90:90:#3);
  \draw[line, fermion] (#1,{#2+#3}) arc (90:270:#3);
  \node[left=3pt] at ({#1-#3},#2) {$#4$};
  \vertexdefaut{#1,{#2+#3}}
}

% Extrémités du propagateur : #1 x gauche, #2 x droit
\providecommand{\bornes}[2]{%
  \node[below=4pt] at (#1,0) {$\tau$};
  \node[below=4pt] at (#2,0) {$\tau'$};
  \fill (#1,0) circle (0.05);
  \fill (#2,0) circle (0.05);
}

\begin{tikzpicture}[
    line/.style={thick},
    >=Stealth,
    fermion/.style={
      postaction={decorate},
      decoration={markings, mark=at position 0.55 with {\arrow{Stealth}}}
    },
]

\def\colL{0}
\def\colR{7.8}
\def\rowA{0}
\def\rowB{-3.9}
\def\rowC{-7.8}

% ===================== (a) double diffusion : tau_2 puis tau_1 =====================
\begin{scope}[shift={(\colL,\rowA)}]
  \draw[line, fermion] (0,0)   -- (1.4,0);
  \draw[line, fermion] (1.8,0) -- (3.0,0);
  \draw[line, fermion] (3.4,0) -- (4.8,0);
  \node[above=3pt] at (0.7,0) {$n$};
  \node[above=3pt] at (2.4,0) {$n_1$};
  \node[above=3pt] at (4.1,0) {$m$};

  \vertexdefaut{1.6,0}
  \node[above=9pt] at (1.6,0) {$\tau_2$};
  \node[below=6pt] at (1.6,0) {$M_{n_1n}$};

  \vertexdefaut{3.2,0}
  \node[above=9pt] at (3.2,0) {$\tau_1$};
  \node[below=6pt] at (3.2,0) {$M_{mn_1}$};

  \bornes{0}{4.8}
  \node[font=\small\bfseries] at (-0.7,2.2) {(a)};
\end{scope}

% ===================== (b) une diffusion (tau_1) + boucle (tau_2) =====================
\begin{scope}[shift={(\colR,\rowA)}]
  \draw[line, fermion] (0,0)   -- (2.2,0);
  \draw[line, fermion] (2.6,0) -- (4.8,0);
  \node[above=3pt] at (1.1,0) {$n$};
  \node[above=3pt] at (3.7,0) {$m$};

  \vertexdefaut{2.4,0}
  \node[below=6pt] at (2.4,0) {$\tau_1$, $M_{mn}$};

  \boucledefaut{2.4}{1.35}{0.55}{n_2}
  \node[above=6pt] at (2.4,1.90) {$\tau_2$, $M_{n_2n_2}$};

  \bornes{0}{4.8}
  \node[font=\small\bfseries] at (-0.7,2.2) {(b)};
\end{scope}

% ===================== (c) une diffusion (tau_2) + boucle (tau_1) =====================
\begin{scope}[shift={(\colL,\rowB)}]
  \draw[line, fermion] (0,0)   -- (2.2,0);
  \draw[line, fermion] (2.6,0) -- (4.8,0);
  \node[above=3pt] at (1.1,0) {$n$};
  \node[above=3pt] at (3.7,0) {$m$};

  \vertexdefaut{2.4,0}
  \node[below=6pt] at (2.4,0) {$\tau_2$, $M_{mn}$};

  \boucledefaut{2.4}{1.35}{0.55}{n_1}
  \node[above=6pt] at (2.4,1.90) {$\tau_1$, $M_{n_1n_1}$};

  \bornes{0}{4.8}
  \node[font=\small\bfseries] at (-0.7,2.2) {(c)};
\end{scope}

% ===================== (d) double diffusion : tau_1 puis tau_2 =====================
\begin{scope}[shift={(\colR,\rowB)}]
  \draw[line, fermion] (0,0)   -- (1.4,0);
  \draw[line, fermion] (1.8,0) -- (3.0,0);
  \draw[line, fermion] (3.4,0) -- (4.8,0);
  \node[above=3pt] at (0.7,0) {$n$};
  \node[above=3pt] at (2.4,0) {$n_2$};
  \node[above=3pt] at (4.1,0) {$m$};

  \vertexdefaut{1.6,0}
  \node[above=9pt] at (1.6,0) {$\tau_1$};
  \node[below=6pt] at (1.6,0) {$M_{n_2n}$};

  \vertexdefaut{3.2,0}
  \node[above=9pt] at (3.2,0) {$\tau_2$};
  \node[below=6pt] at (3.2,0) {$M_{mn_2}$};

  \bornes{0}{4.8}
  \node[font=\small\bfseries] at (-0.7,2.2) {(d)};
\end{scope}

% ===================== (e) propagation libre + deux boucles =====================
\begin{scope}[shift={(\colL,\rowC)}]
  \draw[line, fermion] (0,0) -- (4.8,0);
  \node[below=3pt] at (2.4,0) {$n$};

  \boucledefaut{1.3}{1.35}{0.55}{n_1}
  \node[above=6pt] at (1.3,1.90) {$\tau_1$, $M_{n_1n_1}$};

  \boucledefaut{3.5}{1.35}{0.55}{n_2}
  \node[above=6pt] at (3.5,1.90) {$\tau_2$, $M_{n_2n_2}$};

  \bornes{0}{4.8}
  \node[font=\small\bfseries] at (-0.7,2.2) {(e)};
\end{scope}

% ===================== (f) propagation libre + boucle à deux vertex =====================
\begin{scope}[shift={(\colR,\rowC)}]
  \draw[line, fermion] (0,0) -- (4.8,0);
  \node[below=3pt] at (2.4,0) {$n$};

  % boucle : centre (2.4,1.35), demi-axes 1.1 et 0.55
  \draw[line, fermion] (1.3,1.35) arc (180:0:1.1 and 0.55);
  \draw[line, fermion] (3.5,1.35) arc (0:-180:1.1 and 0.55);
  \node[above=2pt] at (2.4,1.90) {$n_2$};
  \node[below=2pt] at (2.4,0.80) {$n_1$};

  \vertexdefaut{1.3,1.35}
  \node[left=4pt] at (1.3,1.35) {$\tau_1$};
  \vertexdefaut{3.5,1.35}
  \node[right=4pt] at (3.5,1.35) {$\tau_2$};

  \node[font=\small\bfseries] at (-0.7,2.2) {(f)};
  \bornes{0}{4.8}
\end{scope}

\end{tikzpicture}

    \caption[Diagrammes de Feynman des termes du deuxième ordre de la série de Dyson.]{Diagrammes de Feynman des termes du deuxième ordre de la série de Dyson, dans l'ordre des six appariements de Wick. (a) et (d) sont les deux diagrammes connectés, qui diffèrent par l'ordre temporel des deux diffusions sur le défaut. (b), (c), (e) et (f) contiennent au moins une boucle fermée déconnectée du propagateur principal et sont annulés par le dénominateur de l'équation~\eqref{eqn:G}.}
    \label{fig:dyson_ordre2}
\end{figure}
\noindent
De ces six diagrammes, seuls les diagrammes (a) et (d) sont connectés. Ils diffèrent par l'ordre temporel des deux diffusions sur le défaut. Cet ordre est toutefois arbitraire, puisque les temps $\tau_1$ et $\tau_2$ sont intégrés : les deux diagrammes sont donc égaux et leur somme annule le $2!$ au dénominateur du préfacteur des termes de deuxième ordre.

En théorie, nous devrions faire cet exercice pour tous les ordres. En pratique, les calculs deviennent lourds très rapidement, comme l'illustre déjà le deuxième ordre. Au lieu de nous perdre dans des calculs interminables, nous pouvons gagner du temps en reconnaissant une régularité : à l'ordre $N$, il y a $(N+1)!$ appariements et seuls $N!$ d'entre eux sont connectés. Qui plus est, ces $N!$ diagrammes sont tous équivalents et annulent exactement le $N!$ au dénominateur du préfacteur de chaque ordre. Ainsi, en ne considérant que les diagrammes connectés, nous obtenons une série perturbative pour la fonction de Green de Matsubara du système perturbé, exprimée en termes de diagrammes de Feynman et illustrée à la figure~\ref{fig:dyson_serie}.
\begin{figure}[h]
    \centering
    % Développement complet de la série de Dyson pour la fonction de Green
% électronique en présence d'un seul défaut.
% Style identique aux diagrammes du chapitre Raman (tikz/e1.tex, tikz/pd.tex).
% Compatible pdflatex. Requiert : tikz, decorations.markings, arrows.meta, calc.
% Géométrie : segments de 0.58, vertex de rayon 0.2, séparateurs de 0.7.
% Ajustée pour tenir dans la justification du mémoire (470 pt).

% Vertex électron-défaut : cercle blanc + croix (cf. tikz/règles_Raman.tex)
\providecommand{\vertexdefaut}[1]{%
  \draw[fill=white] (#1) circle (0.2);
  \draw[line width=1pt] ($(#1)+(-0.14,-0.14)$) -- ($(#1)+(0.14,0.14)$);
  \draw[line width=1pt] ($(#1)+(-0.14,0.14)$) -- ($(#1)+(0.14,-0.14)$);
}

% Extrémités d'un terme : #1 x gauche, #2 x droit
\providecommand{\bornesserie}[2]{%
  \fill (#1,0) circle (0.05);
  \fill (#2,0) circle (0.05);
  \node[below=2pt] at (#1,0) {\small$\tau$};
  \node[below=2pt] at (#2,0) {\small$\tau'$};
}

\begin{tikzpicture}[
    line/.style={thick},
    >=Stealth,
    fermion/.style={
      postaction={decorate},
      decoration={markings, mark=at position 0.55 with {\arrow{Stealth}}}
    },
]

\node[anchor=east] at (1.55,0) {$\mathcal{G}_{mn}(\tau'-\tau) =$};

% ---- ordre 0 : propagation libre ----
\draw[line, fermion] (1.75,0) -- (2.90,0);
\bornesserie{1.75}{2.90}

\node at (3.25,0) {$+$};

% ---- ordre 1 : un vertex ----
\draw[line, fermion] (3.60,0) -- (4.18,0);
\draw[line, fermion] (4.58,0) -- (5.16,0);
\vertexdefaut{4.38,0}
\bornesserie{3.60}{5.16}

\node at (5.51,0) {$+$};

% ---- ordre 2 : deux vertex ----
\draw[line, fermion] (5.86,0) -- (6.44,0);
\draw[line, fermion] (6.84,0) -- (7.42,0);
\draw[line, fermion] (7.82,0) -- (8.40,0);
\vertexdefaut{6.64,0}
\vertexdefaut{7.62,0}
\bornesserie{5.86}{8.40}

\node at (8.75,0) {$+$};

% ---- ordre 3 : trois vertex ----
\draw[line, fermion] (9.10,0)  -- (9.68,0);
\draw[line, fermion] (10.08,0) -- (10.66,0);
\draw[line, fermion] (11.06,0) -- (11.64,0);
\draw[line, fermion] (12.04,0) -- (12.62,0);
\vertexdefaut{9.88,0}
\vertexdefaut{10.86,0}
\vertexdefaut{11.84,0}
\bornesserie{9.10}{12.62}

\node[anchor=west] at (12.95,0) {$+\ \cdots$};

\end{tikzpicture}

    \caption[Équation de Dyson pour la fonction de Green de Matsubara électronique en présence d'un défaut.]{Équation de Dyson pour la fonction de Green de Matsubara électronique en présence d'un défaut.}
    \label{fig:dyson_serie}
\end{figure}

\noindent
Cette série de diagrammes est l'équation de Dyson pour la diffusion sur un seul défaut et peut être sommée exactement, ce qui n'est pas surprenant car il s'agit bien d'un problème de diffusion à un corps. Généralement, dans le contexte de la diffusion sur des défauts, l'équation de Dyson est résolue en introduisant la matrice $T$ définie diagrammatiquement à la figure~\ref{fig:matrice_T}.
\begin{figure}[h]
    \centering
    % Définition diagrammatique de la matrice T : développement complet (première
% ligne) et forme récursive de Dyson (deuxième ligne).
% Style identique aux diagrammes du chapitre Raman (tikz/e1.tex, tikz/pd.tex).
% Compatible pdflatex. Requiert : tikz, decorations.markings, arrows.meta, calc.
% Le blob hachuré est tracé à la main (clip + traits) pour éviter d'avoir à
% charger la bibliothèque patterns dans le préambule.

% Vertex électron-défaut : cercle blanc + croix (cf. tikz/règles_Raman.tex)
\providecommand{\vertexdefaut}[1]{%
  \draw[fill=white] (#1) circle (0.2);
  \draw[line width=1pt] ($(#1)+(-0.14,-0.14)$) -- ($(#1)+(0.14,0.14)$);
  \draw[line width=1pt] ($(#1)+(-0.14,0.14)$) -- ($(#1)+(0.14,-0.14)$);
}

% Matrice T : cercle hachuré. #1 = coordonnée du centre.
\providecommand{\blobT}[1]{%
  \begin{scope}
    \clip (#1) circle (0.28);
    \foreach \i in {-8,...,8} {
      \draw[line width=0.4pt] ($(#1)+(0.08*\i-0.45,-0.45)$) -- ($(#1)+(0.08*\i+0.45,0.45)$);
    }
  \end{scope}
  \draw[thick, fill=none] (#1) circle (0.28);
}

\begin{tikzpicture}[
    line/.style={thick},
    >=Stealth,
    fermion/.style={
      postaction={decorate},
      decoration={markings, mark=at position 0.55 with {\arrow{Stealth}}}
    },
]

% ============ première ligne : développement complet ============
\node[anchor=east] at (1.5,0) {$T_{mn} =$};
\blobT{1.9,0}
\node at (2.6,0) {$=$};

% un vertex
\vertexdefaut{3.3,0}
\node at (3.9,0) {$+$};

% deux vertex
\vertexdefaut{4.5,0}
\draw[line, fermion] (4.7,0) -- (5.7,0);
\vertexdefaut{5.9,0}
\node at (6.5,0) {$+$};

% trois vertex
\vertexdefaut{7.1,0}
\draw[line, fermion] (7.3,0) -- (8.3,0);
\vertexdefaut{8.5,0}
\draw[line, fermion] (8.7,0) -- (9.7,0);
\vertexdefaut{9.9,0}
\node[anchor=west] at (10.35,0) {$+\ \cdots$};

% ============ deuxième ligne : forme récursive ============
\node at (2.6,-1.7) {$=$};

\vertexdefaut{3.3,-1.7}
\node at (3.9,-1.7) {$+$};

\vertexdefaut{4.5,-1.7}
\draw[line, fermion] (4.7,-1.7) -- (5.72,-1.7);
\blobT{6.0,-1.7}

\end{tikzpicture}

    \caption[Expression diagrammatique de la matrice $T$ pour la diffusion sur un seul défaut.]{Expression diagrammatique de la matrice $T$ pour la diffusion sur un seul défaut. La première ligne donne le développement complet, la deuxième la forme récursive.}
    \label{fig:matrice_T}
\end{figure}

\noindent
La matrice $T$ pour un seul défaut joue essentiellement le même rôle que la \textit{self}-énergie électronique dans le contexte de l'interaction électron-électron ou électron-phonon. En traduisant la première ligne de la figure~\ref{fig:matrice_T} à l'aide des règles de Feynman, nous obtenons la série
\begin{equation}
    T = M + M\,\mathcal{G}^{(0)}M + M\,\mathcal{G}^{(0)}M\,\mathcal{G}^{(0)}M+\dots,
    \label{eqn:T_serie}
\end{equation}
où les produits sont matriciels sur les indices de bande et convolutifs sur le temps imaginaire. La deuxième ligne de la figure~\ref{fig:matrice_T} exprime la même série sous forme récursive
\begin{equation}
    T = M + M\,\mathcal{G}^{(0)}T,
    \label{eqn:T_recursive}
\end{equation}
qui se résout exactement. Comme le défaut est statique, $M$ ne dépend pas de la fréquence et la matrice $T$ est diagonale en fréquence de Matsubara, où les convolutions temporelles deviennent de simples produits
\begin{equation}
    T(i\omega_j) = \left[1-M\,\mathcal{G}^{(0)}(i\omega_j)\right]^{-1}M,
    \label{eqn:T_fermee}
\end{equation}
ou, explicitement en indices de bande,
\begin{equation}
    T_{mn}(i\omega_j) = M_{mn}+\sum_{n_1}M_{mn_1}\mathcal{G}^{(0)}_{n_1}(i\omega_j)T_{n_1n}(i\omega_j),\qquad \mathcal{G}^{(0)}_n(i\omega_j)=\frac{1}{i\hbar\omega_j-\xi_n},
    \label{eqn:T_indices}
\end{equation}
avec $\xi_n = \varepsilon_n-\mu$, où $\mu$ est le potentiel chimique. C'est là tout l'intérêt de la matrice $T$ : une seule inversion de matrice, de la taille du nombre de bandes retenues, resomme la série à tous les ordres, sans hypothèse sur la force du potentiel de défaut. Cette resommation à tous les ordres n'est pas un simple raffinement formel : elle est indispensable pour les défauts considérés dans ce mémoire. Le potentiel d'une lacune correspond à la suppression d'un atome du réseau, une perturbation dont les éléments de matrice sont comparables à la largeur de bande — il n'existe aucun petit paramètre justifiant la troncature de la série de Dyson à un ordre fini. La matrice $T$, exacte pour un centre diffuseur unique, demeure valide dans cette limite de couplage fort (dite limite unitaire). 

Il est possible de réécrire l'équation de Dyson en termes de la matrice $T$. Cela est illustré diagrammatiquement à la figure~\ref{fig:dyson_matrice_T}.
\begin{figure}[h]
    \centering
    % Série de Dyson resommée en termes de la matrice T :
%   G_{mn}(tau'-tau) = G^(0) + G^(0) T G^(0)
% Style identique aux diagrammes du chapitre Raman (tikz/e1.tex, tikz/pd.tex).
% Compatible pdflatex. Requiert : tikz, decorations.markings, arrows.meta, calc.
%
% Convention temporelle : temps croissant de gauche à droite, donc
% tau < tau_1 < tau_2 < tau'. Le deuxième terme se lit
%   G^(0)_m(tau'-tau_2) T_{mn}(tau_2-tau_1) G^(0)_n(tau_1-tau).
% Si l'équation du texte adopte l'ordre inverse, il suffit d'échanger les
% étiquettes tau_1 et tau_2 ci-dessous.

% Vertex électron-défaut : cercle blanc + croix (cf. tikz/règles_Raman.tex)
\providecommand{\vertexdefaut}[1]{%
  \draw[fill=white] (#1) circle (0.2);
  \draw[line width=1pt] ($(#1)+(-0.14,-0.14)$) -- ($(#1)+(0.14,0.14)$);
  \draw[line width=1pt] ($(#1)+(-0.14,0.14)$) -- ($(#1)+(0.14,-0.14)$);
}

% Matrice T : cercle hachuré. Définition identique à celle de tikz/matrice_T.tex.
\providecommand{\blobT}[1]{%
  \begin{scope}
    \clip (#1) circle (0.28);
    \foreach \i in {-8,...,8} {
      \draw[line width=0.4pt] ($(#1)+(0.08*\i-0.45,-0.45)$) -- ($(#1)+(0.08*\i+0.45,0.45)$);
    }
  \end{scope}
  \draw[thick, fill=none] (#1) circle (0.28);
}

\begin{tikzpicture}[
    line/.style={thick},
    >=Stealth,
    fermion/.style={
      postaction={decorate},
      decoration={markings, mark=at position 0.55 with {\arrow{Stealth}}}
    },
]

\node[anchor=east] at (1.55,0) {$\mathcal{G}_{mn}(\tau'-\tau) =$};

% ---- premier terme : propagateur libre ----
\draw[line, fermion] (1.75,0) -- (3.35,0);
\node[above=3pt] at (2.55,0) {$n$};
\fill (1.75,0) circle (0.05);
\fill (3.35,0) circle (0.05);
\node[below=2pt] at (1.75,0) {\small$\tau$};
\node[below=2pt] at (3.35,0) {\small$\tau'$};

\node at (3.85,0) {$+$};

% ---- deuxième terme : G^(0) T G^(0) ----
\draw[line, fermion] (4.35,0) -- (5.95,0);
\node[above=3pt] at (5.15,0) {$n$};
\blobT{6.23,0}
\node[above=8pt] at (6.23,0) {$T_{mn}$};
\draw[line, fermion] (6.51,0) -- (8.11,0);
\node[above=3pt] at (7.31,0) {$m$};

\fill (4.35,0) circle (0.05);
\fill (8.11,0) circle (0.05);
\node[below=2pt] at (4.35,0) {\small$\tau$};
% étiquettes splayées en biais : les deux points d'attache du blob ne sont
% distants que de 0.56, insuffisant pour deux étiquettes centrées
\node[anchor=north east] at (5.99,-0.10) {\small$\tau_1$};
\node[anchor=north west] at (6.47,-0.10) {\small$\tau_2$};
\node[below=2pt] at (8.11,0) {\small$\tau'$};

\end{tikzpicture}

    \caption[Série de Dyson pour la fonction de Green de Matsubara électronique en présence d'un seul défaut exprimée en termes de la matrice $T$.]{Série de Dyson pour la fonction de Green de Matsubara électronique en présence d'un seul défaut, exprimée en termes de la matrice $T$. Les temps $\tau_1$ et $\tau_2$ sont intégrés.}
    \label{fig:dyson_matrice_T}
\end{figure}

\noindent
Finalement, nous obtenons une expression mathématique pour la fonction de Green de Matsubara du système perturbé en traduisant la série de diagrammes de la figure~\ref{fig:dyson_matrice_T} à l'aide des règles de Feynman. Cela donne
\begin{equation}
    \mathcal{G}_{mn}(\tau'-\tau) = \delta_{mn}\mathcal{G}^{(0)}_n(\tau'-\tau) + \int_0^\beta d\tau_1\,\int_0^\beta d\tau_2\,\mathcal{G}^{(0)}_m(\tau'-\tau_2)T_{mn}(\tau_2-\tau_1)\mathcal{G}^{(0)}_n(\tau_1-\tau).
\end{equation}

\section{Généralisation à des systèmes désordonnés}
